%! @@TITLE@@ — Unit @@UNITINT@@, Lesson @@LESSONID@@ — Guided Notes & Practice (Answer Key)
% MIRROR of ../notes/main.tex (Units 3+ packet-first, LESSON_SHAPE.md §0a). Copy the blank
% VERBATIM, then change exactly these:
%   \usepackage{@@PREFIX@@-boxes}      ->  \usepackage{@@PREFIX@@-key}   (pulls in -boxes; never load both)
%   the \pageheader second argument    ->  Guided Notes \& Practice --- Answer Key
%   \blank{W}                          ->  \ans{short answer}           (no wider than the blank)
%   \writelines{n}                     ->  exactly n \ansline{...}      (each under ~95 characters)
% Every \begin{work} block stays BYTE-IDENTICAL — the blank and key print the same solution
% lines (the work environment shows them only in key mode), so heights cannot drift.
% \ans{} is TEXT MODE: never inside $...$; wrap \sqrt/\dfrac/^/_ in $...$ inside it.
% NO teachernote here — teacher prose goes in the lesson plan.
% Prove it: diff the two files and only the lines above may differ; then the page counts match.
\documentclass[12pt]{article}
\usepackage{@@PREFIX@@-article}
\usepackage{@@PREFIX@@-key}
\newcommand{\rung}[1]{{\color{goldacc}\ifcase#1\or$\star$\or$\star\star$\or$\star\star\star$\fi}}

\begin{document}

\pageheader{Unit @@UNITINT@@, Lesson @@LESSONID@@}{Guided Notes \& Practice --- Answer Key}

\vspace{0.10in}

\boxguard[20]
\begin{notesbox}{Round 1 --- TODO: the first skill}
\small
\textbf{TODO term.} TODO: the definition as one complete printed sentence.

\vspace{6pt}
\textbf{Example.} TODO: the problem.
\begin{work}
  % TODO — identical in the blank.
  x &= x
\end{work}

\vspace{6pt}
\textbf{Your turn.} \textit{Start at \rung{1} and climb.}

\rung{1}\; TODO
\begin{work}
  x &= x
\end{work}

\rung{2}\; TODO
\begin{work}
  x &= x
\end{work}

\rung{3}\; TODO
\begin{work}
  x &= x
\end{work}

\textbf{Show me.} My answer: \ans{TODO}
\end{notesbox}

% TODO: Rounds 2 and 3, mirrored.

\boxguard[20]
\begin{notesbox}{Closing Practice --- score yourself out of 4}
\small
\begin{enumerate}[leftmargin=1.9em, itemsep=8pt]
  \item TODO (Round 1) \ans{TODO}
  \item TODO (Round 2) \ans{TODO}
  \item TODO (Round 3) \ans{TODO}
  \item TODO (mixed) \ans{TODO}
\end{enumerate}
\textbf{Bonus (+1).} TODO: a student-life context. \ans{TODO}
\end{notesbox}

\end{document}
